\section{Product Overview}

\subsection{Product Perspective}

The Medication Safety Assistant is a standalone responsive Web prototype that
coordinates external data and processing services. It is not an electronic
medical record, a prescribing system, a diagnostic device, or an autonomous
clinical decision-maker.

At the product boundary, patients and caregivers provide prescription images,
manual text, typed requests, voice input, relevant health information, schedule
confirmation, adherence status, and possible adverse-reaction reports.
Healthcare professionals interact with a simulated review view. The browser
provides camera, microphone, and notification capabilities when permission is
granted.

The high-level information flow is:

\begin{enumerate}
  \item The patient provides medication information and relevant health
  context through an image, voice, manual entry, or typed request.
  \item The system extracts editable fields and marks low-confidence, missing,
  or conflicting information. Unreliable extraction returns to re-entry or
  correction rather than continuing to assessment.
  \item The patient verifies that the extracted details match the source. This
  confirmation concerns transcription accuracy, not medical correctness. Any
  edit invalidates the previous normalized result and starts normalization again.
  \item The system performs NLP, DAV identification, ingredient normalization,
  and reviewed DDInter mapping, then checks whether the required safety context
  is sufficient. Missing information is requested before reassessment.
  \item The system assesses supported risks and chooses one of three routing
  outcomes: routine/no significant risk detected from available information,
  professional review required or uncertain, or emergency/red flags.
  \item Routine cases receive qualified guidance. Review-required cases remain
  pending until reviewed; if review is unavailable, the system gives
  conservative guidance and recommends contacting a qualified doctor or
  pharmacist. Emergency cases receive immediate urgent guidance without waiting
  for professional review in the prototype.
  \item After reviewing and confirming the final medication plan, the patient
  may explicitly opt in to reminders. Dose responses and possible adverse
  reactions are then recorded for follow-up and triage.
\end{enumerate}

OCR transforms printed or digital Vietnamese prescriptions into editable
fields. Speech-to-text transforms voice into an editable transcript. NLP
identifies one or more intents and medication entities from verified text and
may explain structured results, but it does not independently infer safety.
DAV publication data supports Vietnamese product and active-ingredient
identification. DDInter 2.0 supplies supported interaction information after a
reliable mapping is established. External service failure, ambiguous mapping,
or insufficient data produces an unavailable, unresolved, incomplete, or
unassessed state rather than a safe result. A simulated professional reviewer
may approve guidance, request more information, or leave a case pending. An
unreviewed case can never be represented as professionally approved.

\subsection{Assumptions and Dependencies}

\begin{longtable}{|L{0.16\textwidth}|L{0.43\textwidth}|L{0.31\textwidth}|}
\hline
\textbf{Type} & \textbf{Assumption / dependency} & \textbf{Effect if invalid or unavailable}\\ \hline
Assumption & Users can read, correct, and explicitly confirm OCR results, voice transcripts, medicines, relevant health context, and schedules. & Unconfirmed information cannot enter safety assessment or create an active reminder.\\ \hline
Dependency & DAV medicine-registration data remains accessible and usable for identifying Vietnamese products and active ingredients. & Identification coverage is reduced; affected products remain unresolved. The system must show the source date and limitations.\\ \hline
Dependency & Prepared DAV-derived normalization data and DAV-to-DDInter mappings can be versioned and reviewed before use. & Unreviewed or ambiguous mappings must be excluded from automatic assessment until corrected.\\ \hline
Dependency & DDInter 2.0 remains accessible and its permitted use supports the course prototype. & Interaction checks may use a permitted controlled snapshot or be unavailable; the interface must disclose coverage and date.\\ \hline
Dependency & OCR can process sufficiently clear printed or digital Vietnamese prescriptions. & Poor-quality, complex, or handwritten content requires manual correction or entry and remains unsupported when unresolved.\\ \hline
Dependency & Speech-to-text supports the user's language, microphone, and browser with acceptable recognition. & Users must receive a clear error and retain typed/manual input; an unconfirmed transcript cannot be assessed.\\ \hline
Dependency & The selected NLP service or model can identify intents/entities and explain only supplied structured facts. & NLP output is withheld or replaced with deterministic text when constraints cannot be enforced or the service is unavailable.\\ \hline
Assumption & Professional review is simulated for the course prototype. & No claim may be made that a real clinician has reviewed a case. Pending or unavailable review remains explicit, while emergency and other appropriate cases direct the user to actual professional care.\\ \hline
Assumption & Modern desktop and mobile browsers provide the required Web, camera, microphone, and notification capabilities. & Unsupported permissions or capabilities require graceful fallback to manual entry and in-application schedules.\\ \hline
Constraint & The team uses synthetic or de-identified development and demonstration data. & Real patient deployment, retention, and clinical integration remain outside the initial scope.\\ \hline
\end{longtable}

